% AutoPatch results section draft
% Model: Claude 3.5 Sonnet | Dataset: ARVO-Lite (15 cases)
% Generated by result-analyzer-and-table-gen skill v1.0.0 on 2026-04-02
% TODO: update cross-refs once companion sections exist

\section{Results}
\label{sec:results}

\subsection{Overall Performance}
\label{sec:results-overall}

\autoref{tab:results-main} summarises the performance of Claude~3.5 Sonnet on the
\textsc{ARVO-Lite} benchmark (15~cases).
The model passes quality-assurance container checks on 13 of 15 cases (86.7\%),
demonstrating that the patching harness can reliably set up reproducible
build-and-fuzz environments for almost all targets.
Of those 13~QA-passed cases, the model produces a syntactically valid patched function in
12~instances (patch generation rate 92.3\%), confirming that the LLM rarely
fails to emit a well-formed patch once the environment is ready.
At the end-to-end level, 9~of 15 cases have their crash repaired as verified by
the fuzzing oracle (Tier~1 pass rate 60.0\%), and 8~cases additionally pass
differential debugging — confirming that the fix does not introduce new
regressions (correct patch rate 53.3\%).
These results represent a competitive single-model baseline for autonomous
vulnerability patching on real-world C/C++ targets.
% TODO: update cross-ref once section exists: \autoref{sec:method}

\subsection{Per-Tier Analysis}
\label{sec:results-tiers}

\paragraph{Tier~1 (crash fix).}
The Tier~1 oracle re-runs the original crash-triggering input through the patched
binary under fuzzing.
The model achieves a 60.0\% Tier~1 pass rate (9/15).
As shown in \autoref{fig:results-status}, the majority of cases that do not reach
Tier~1 either stall at the build stage (2~cases: \textsc{PATCH\_BUILD\_SUCCESSFUL})
or fail before producing a formatted patch (1~case: \textsc{PATCH\_FORMAT\_CORRECT}).
Three cases fail outright (\textsc{FAILED}), and one produces a patch that fixes
the crash but is flagged by a fuzzing decode error, which may indicate a harness
artefact rather than a patch correctness problem.

\paragraph{Tier~1 $\rightarrow$ Correct (differential debugging).}
Of the 9~cases that pass Tier~1, 8 additionally survive differential debugging,
yielding a correct patch rate of 53.3\%.
The single regression case highlights that surface-level crash repair does not
guarantee semantic equivalence with the unpatched code path.

\paragraph{Multi-tier fields (Tier~2--4).}
This run did not use the \texttt{--multi-tier} flag; accordingly, \texttt{tier2\_tests\_passed},
\texttt{tier3\_coverage\_pct}, and \texttt{tier4\_diff\_passed} fields are absent.
Results for the extended evaluation pipeline will be reported in a subsequent experiment.

\subsection{Error Analysis}
\label{sec:results-errors}

\autoref{fig:results-status} shows the distribution of
\texttt{max\_patch\_generation\_status} values across all 15~cases.
The most common terminal status is \textsc{PATCH\_PASSES\_CHECKS} (8~cases, 53.3\%),
indicating that patches pass post-generation quality checks and proceed to fuzzing
evaluation — making the fuzzing oracle the primary discriminator between success
and failure.
Three cases (20.0\%) fail completely before generating any patch, most likely
due to build-environment idiosyncrasies that the model cannot resolve within the
allotted build-iteration budget.

\autoref{fig:results-iters} shows the distribution of \texttt{build\_iters} and
\texttt{fix\_crash\_iters} per case.
The average number of build-fix cycles is 2.1, and the average number of
crash-fix iterations is 1.8, both within the budget limits configured for this
run.
Cases that exhaust their iteration budget without a passing build account for the
two \textsc{PATCH\_BUILD\_SUCCESSFUL}-terminal cases.
Average patch generation wall-clock time is 4.7~minutes per case.

\subsection{Ablation Study}
\label{sec:results-ablation}

\textit{[To be filled in after ablation runs.]}
% TODO: add ablation results once runs complete

\subsection{Comparison to Baselines}
\label{sec:results-comparison}

% TODO: update cross-ref once section exists: \autoref{tab:results-comparison}
A comparison table across multiple models (\texttt{paper/tables/results\_comparison.tex})
will be populated once baseline runs (GPT-4o, Gemini~1.5~Pro) are complete.
The schema example value for GPT-4o (\texttt{shr\_correct\_patches}~=~0.583) places
Claude~3.5 Sonnet approximately 5 percentage points lower on this 15-case dataset.
Given the small sample size ($n$~=~15), a difference of one case (6.7~pp) is
within sampling noise; further evaluation on the full \textsc{ARVO} dataset is
needed to establish significance.
